% Isolated research note.  No evenness or full-monodromy conclusion is claimed.
\subsection*{The exact boundary term in a proposed main-component parity argument}

Use the actual square-root convention of (DE-LP2), and suppose that
the reduced polar \(Q=(q=0)\subset\mathbf P^{N-1}_y\) is integral
and meets the open on which \(B(y^{[2]})\) has rank \(N-1\).
Write \(e=\deg\sigma\), so \(\deg Q=2e+1\), and let \(Q^\circ\)
be that maximal-rank open.  These hypotheses are precisely the
proved inputs to the unique main Lang component.  The following
calculation isolates a further divisor which is not yet controlled.

\paragraph{A flag graph, without resolution of singularities.}
For \([y]\in Q^\circ\), let
\[
 R_y=k\sigma(y)^{[2]}\subset W,\qquad
 E_y=\{s\in W:B(y^{[2]})s\in ky\}.
\]
The first is a line and the second a plane containing it.  Form the
closure of the graph of the rational maps
\[
 Q^\circ\longrightarrow
 \mathbf P^{N-1}_{\sigma}\times\operatorname{Flag}(1,2;W),
 \qquad [y]\longmapsto([\sigma(y)],R_y\subset E_y),
\]
and, if desired, normalize this integral projective graph.  Call
the result \(X\), with its projective birational map to \(Q\).
This construction uses no unresolved general resolution theorem
in positive characteristic.

Let \(H\) denote the pullback of the hyperplane line bundle from
\(Q\), and \(L\) the pullback of \(\mathcal O(1)\) from the
\(\sigma\)-factor.  The tautological plane bundle \(\mathcal E\)
is a rank-two subbundle of \(W\otimes\mathcal O_X\).
The closed flag incidence and the coordinatewise square map give
an exact sequence of vector bundles
\[
 0\longrightarrow\mathcal O_X(-2L)
 \longrightarrow\mathcal E
 \longrightarrow\mathcal M\longrightarrow0.                    \tag{MP1}
\]
The last term is a line bundle.  The actual equation
\(B(y^{[2]})s\in ky\), valid densely, extends over the integral
graph as an identity of bundle morphisms.  Since the matrix entries
have degree two in \(y\), it supplies
\[
       \mathcal E\longrightarrow\mathcal O_X(H).
\]
It kills \(\mathcal O_X(-2L)\), and hence induces a nonzero map of
line bundles \(\mathcal M\to\mathcal O_X(H)\).  Let \(D\) be its
effective Cartier zero divisor.  Thus the exact global quotient is
\[
                    \mathcal M=\mathcal O_X(H-D).               \tag{MP2}
\]
The divisor \(D\) is supported over the boundary \(X\setminus
Q^\circ\).  Replacing it by zero, or by an even divisor without
proof, changes the degree calculation.

\paragraph{The resulting degree identity.}
Use the convention that \(\mathbf P_X(\mathcal E)\) parametrizes
lines in \(\mathcal E\), and let
\(\xi=c_1(\mathcal O_{\mathbf P(\mathcal E)}(1))\) be the dual
tautological class.  The inclusion \(\mathcal E\subset
W\otimes\mathcal O_X\) defines a morphism
\(\psi:\mathbf P_X(\mathcal E)\to\mathbf P(W)\), with pullback
hyperplane class \(\xi\).  On the dense maximal-rank open, lines
with \(B(y^{[2]})s\ne0\) correspond to the actual nonzero Lang
solutions by the unique rescaling that makes the proportionality
constant one.  The radical section has smaller dimension.
Consequently \(\psi\) is generically finite of the same degree
\(\delta\) as the main Lang cover.

Put \(n=\dim X=N-2\).  The projective-bundle formula and
multiplicativity of Chern classes in (MP1) give
\[
 \begin{aligned}
 \delta&=\int_{\mathbf P(\mathcal E)}\xi^{n+1}
      =\int_X s_n(\mathcal E)\\
      &=\int_X\sum_{i=0}^{n}(2L)^i(D-H)^{n-i},\\
 \delta&\equiv\int_X(H-D)^n\pmod2.                              \tag{MP3}
 \end{aligned}
\]
These formulas apply to the integral graph using vector-bundle
Chern operations on Chow groups; smoothness of \(X\) is not assumed.
If the class of \(D\) were proved even, the last expression would
equal \(\int_XH^n=\deg Q=2e+1\) modulo two.  The presently missing
step is precisely such a parity control, or an alternative direct
calculation of the last intersection number.  No oddness of
\(\delta\) is asserted here.

The original Stacks Project statements and full displayed proofs
of Tags \texttt{02TX}, \texttt{02U6}, and Section \texttt{02UF}
were opened and read, together with the convention and definitions
in Section \texttt{02TZ}.  In the Stacks quotient convention the
bundle in the projective formula is \(\mathcal E^\vee\);
this is why (MP3) uses the inverse total Chern class of
\(\mathcal E\).  These sources justify the intersection calculation,
not the unresolved evenness of \(D\).

\paragraph{Squared matrix entries do not force even \(D\).}
Over \(R=k[[t]]\), put
\[
 J=\begin{pmatrix}0&1\\1&0\end{pmatrix},\quad
 A=\operatorname{diag}(tJ,J,0),\quad
 B=A^{[2]}=\operatorname{diag}(t^2J,J,0),\quad
 y=(t,0,1,0,0)^{\mathsf t}.
\]
The primitive radical is the fifth coordinate line.
Over the fraction field, a solution of \(Bv=y\) is
\[
          v=(0,t^{-1},0,1,0)^{\mathsf t}.
\]
The saturated plane in \(R^5\) spanned by the radical and \(v\)
has basis
\[
          e_4,\qquad w=(0,1,0,t,0)^{\mathsf t}.
\]
Here \(Bw=t\,y\).  The quotient-line map to the primitive
\(y\)-line therefore vanishes to order one.  Thus the local
boundary divisor has odd multiplicity despite every matrix
coefficient being a square.  This local example is not claimed
to satisfy the global Cartier or birational-radical hypotheses.
It refutes only the inference from squared coefficients alone.

Equivalently, after an alternating Smith normal form for \(A\),
the elementary valuations of \(B=A^{[2]}\) are even, but the
coordinates of the transformed vector \(U^{[2]}y\) can have odd
valuation.  The least clearing exponent for a solution is
controlled by differences \(2a_i-v((U^{[2]}y)_i)\); there is no
formal parity conclusion for these differences.

\paragraph{A separate strict-transform warning.}
In the exact \(N=5\) countermodel (PH2.9), consider
\(p=[0:0:1:0:0]\).  On \(x_2=1\),
\[
 q=x_0^3x_1^2+x_1^3+x_1^4+x_0^2x_3x_4(x_3+x_4).
\]
The reduced radical ideal has order one at \(p\), whereas
the equation \(q\) has order three.  Blow up \(p\), then
continue by graph modifications resolving the radical map.
At the generic point of the first exceptional divisor the
radical is already nonzero after removing its order-one factor,
so that divisor persists.  Dividing \(q\) by the square of the
radical factor leaves that exceptional divisor with multiplicity
one.  Therefore the expression \(H+2L\) describes a residual
total divisor, not automatically the strict main transform.
This warning is independent of the quotient-line divisor \(D\)
in (MP2).  Neither warning is an actual counterexample to the
desired parity for the geometric Cartier tensor.
